\notechapter{N-20260615}{Surface Integrals, Gauss Theorem, and Stokes Theorem}

\section{Why this note is split this way}

Surface integration is where several notations begin to look dangerously similar.  The guiding principle of this note is that every integral should be identified by the geometric object it integrates over and by the degree of the quantity being integrated.  A scalar surface integral measures density against area; a flux integral measures a two-dimensional oriented flow; Gauss turns a two-dimensional boundary flux into a three-dimensional divergence integral; Stokes turns a one-dimensional boundary circulation into a two-dimensional curl integral.

This is not a taxonomy for its own sake.  It is the first place where the general formula
\[
  \int_{\partial M}\omega=\int_M d\omega
\]
starts to become visible behind vector calculus.



\section{Four similar-looking integrals}

The first task is to separate four objects:
\[
  \iint_S f(x,y,z)\,dS,\qquad
  \iint_S P\,dy\,dz+Q\,dz\,dx+R\,dx\,dy,
\]
\[
  \iiint_\Omega \nabla\cdot F\,dV,\qquad
  \oint_{\partial S}P\,dx+Q\,dy+R\,dz.
\]
The first is a scalar surface integral: it integrates a scalar density over area.  The second is a flux integral: it integrates a vector field through an oriented surface.  The third is the volume integral of divergence.  The fourth is circulation along an oriented boundary curve.

The common mistake is to treat all of them as ``surface integral formulas.''  They are not the same mathematical object.  Their relation is supplied by Gauss' theorem and Stokes' theorem.

\section{First-kind surface integrals}

Let $S$ be a smooth surface parametrized by
\[
  r(u,v)=(x(u,v),y(u,v),z(u,v)),\qquad (u,v)\in D.
\]
Then
\[
  dS=|r_u\times r_v|\,du\,dv,\qquad
  \iint_S f\,dS=\iint_D f(r(u,v))|r_u\times r_v|\,du\,dv.
\]
This is the surface analogue of the first-kind line integral $\int_C f\,ds$.

For a graph $z=z(x,y)$,
\[
  r(x,y)=(x,y,z(x,y)),\qquad
  r_x\times r_y=(-z_x,-z_y,1),
\]
so
\[
  dS=\sqrt{1+z_x^2+z_y^2}\,dx\,dy.
\]
For a graph $x=x(y,z)$ or $y=y(z,x)$, analogous formulas hold:
\[
  dS=\sqrt{1+x_y^2+x_z^2}\,dy\,dz,\qquad
  dS=\sqrt{1+y_z^2+y_x^2}\,dz\,dx.
\]


\begin{figure}[H]
\centering
\includegraphics[width=0.86\textwidth,height=0.74\textheight,keepaspectratio]{figures/notes-md/N-20260615/N-20260615-fig-01-surface-area-element-and-projection.png}
\caption{A surface area element and its projection.  The factor between \(dS\) and the projected area element is determined by the angle between the surface normal and the projection direction.}
\end{figure}

\section{Template surfaces}

The common templates are worth collecting because they prevent repeated derivations.

For a plane $z=ax+by+c$,
\[
  dS=\sqrt{1+a^2+b^2}\,dx\,dy.
\]
For a sphere $x^2+y^2+z^2=R^2$, spherical coordinates give
\[
  dS=R^2\sin\varphi\,d\varphi\,d\theta.
\]
For a cylinder $x=R\cos\theta$, $y=R\sin\theta$, $z=z$,
\[
  dS=R\,d\theta\,dz.
\]
For a cone $z=\sqrt{x^2+y^2}$, use $x=r\cos\theta$, $y=r\sin\theta$, $z=r$, so
\[
  |r_r\times r_\theta|=\sqrt2\,r.
\]


\begin{figure}[H]
\centering
\includegraphics[width=0.70\textwidth,height=0.74\textheight,keepaspectratio]{figures/notes-md/N-20260615/N-20260615-fig-04-torus-example-surface.png}
\caption{A torus as a standard example of a parametrized surface.}
\end{figure}

\section{Example pattern: spherical cap}

For a spherical cap on $x^2+y^2+z^2=R^2$, the scalar integral of $f$ should normally be written in spherical coordinates.  If the cap is cut by $z=h$, then $z=R\cos\varphi$ gives
\[
  0\le\varphi\le \arccos(h/R).
\]
Thus
\[
  \iint_S f\,dS=
  \int_0^{2\pi}\int_0^{\arccos(h/R)}
  f(R\sin\varphi\cos\theta,R\sin\varphi\sin\theta,R\cos\varphi)
  R^2\sin\varphi\,d\varphi\,d\theta.
\]
The important part of the example is not the final antiderivative.  It is recognizing the right parameter domain.

\begin{figure}[H]
\centering
\includegraphics[width=0.78\textwidth,height=0.74\textheight,keepaspectratio]{figures/notes-md/N-20260615/N-20260615-fig-02-spherical-surface-patch-and-projection-domain.png}
\caption{A spherical surface patch and its projection domain.  This is the geometric model behind many spherical-cap surface integrals.}
\end{figure}

\begin{figure}[H]
\centering
\includegraphics[width=0.78\textwidth,height=0.74\textheight,keepaspectratio]{figures/notes-md/N-20260615/N-20260615-fig-03-spherical-surface-with-bounding-cylinder.png}
\caption{A related spherical surface cut out by a vertical boundary.  The dashed curves indicate the hidden projection and bounding geometry.}
\end{figure}





\section{Second-kind surface integrals as flux}

Let $F=(P,Q,R)$.  The flux through an oriented surface $S$ is
\[
  \iint_S F\cdot n\,dS.
\]
In differential notation this is
\[
  \iint_S P\,dy\,dz+Q\,dz\,dx+R\,dx\,dy.
\]
For a parametrization $r(u,v)$,
\[
  \iint_S F\cdot n\,dS
  =\iint_D F(r(u,v))\cdot(r_u\times r_v)\,du\,dv,
\]
where the order of $u,v$ determines orientation.

For a graph $z=z(x,y)$ with upward orientation,
\[
  n\,dS=(-z_x,-z_y,1)\,dx\,dy,
\]
so
\[
  \iint_S F\cdot n\,dS
  =\iint_D[-Pz_x-Qz_y+R]\,dx\,dy.
\]
For downward orientation, change the sign.


\section{Model flux computations: closed, capped, singular, and sheeted}

The useful flux examples are not a random list of surfaces; they represent four different decisions about orientation and theorem choice.

\subsection{Closed sphere: Gauss beats parametrization}

Let \(F(x,y,z)=(x,y,z)\) and let \(S\) be the sphere \(x^2+y^2+z^2=R^2\) with outward orientation.  Directly, the outward unit normal is \(n=x/R\), so
\[
  F\cdot n=R,\qquad dS=R^2\sin\varphi\,d\varphi\,d\theta.
\]
Hence
\[
  \iint_SF\cdot n\,dS=R\cdot4\pi R^2=4\pi R^3.
\]
Gauss gives the same answer faster:
\[
  \nabla\cdot F=3,\qquad
  \iiint_{B_R}3\,dV=3\cdot\frac{4}{3}\pi R^3=4\pi R^3.
\]
The theorem choice is dictated by closedness of the surface.

\subsection{Open surface: close it and subtract the cap}

For the upper hemisphere \(S\) of radius \(R\), outward orientation, Gauss cannot be applied to \(S\) alone because it is not closed.  Close it by adding the disk \(D\) in the plane \(z=0\).  For \(F=(x,y,z)\), the flux through \(D\) is zero because \(F\cdot n=0\) there.  Therefore the hemisphere flux equals the full closed half-ball flux:
\[
  \iint_SF\cdot n\,dS=\iiint_{z\ge0,\ x^2+y^2+z^2\le R^2}3\,dV=2\pi R^3.
\]
For another field, the cap flux might not vanish; subtracting it is the whole point of the method.

\subsection{Cone: orientation is the computation}

For the cone \(z=r\), \(0\le r\le a\), use
\[
  r(r,\theta)=(r\cos\theta,r\sin\theta,r).
\]
Then
\[
  r_r\times r_\theta=(-r\cos\theta,-r\sin\theta,r),
\]
which has positive \(z\)-component.  Thus it is the upward orientation.  If a problem asks for outward orientation of the cone as the side of a solid below its tip, the sign may be the opposite.  Here the formula is easy; choosing the correct normal is the real work.

\subsection{Two sheets over one projection}

For a sphere cut by projection onto the \(xy\)-plane, upper and lower sheets have opposite normal components.  A flux written as
\[
  \iint_S R\,dx\,dy
\]
really means integrating the two-form \(R\,dx\wedge dy\).  On the upper sheet, upward orientation contributes \(+R\,dx\,dy\); on the lower sheet, the same geometric outward orientation may contribute \(-R\,dx\,dy\).  Thus sheets may add or cancel depending on the chosen orientation and the component of the flux form.

This is the differential-form reason that projection formulas are dangerous when treated as scalar area formulas.

\section{Gauss theorem}

\begin{theorem}[Gauss divergence theorem]
For a compact region \(\Omega\) with oriented boundary \(\partial\Omega\),
\[
  \int_{\partial\Omega} F\cdot n\,dS=\int_\Omega \nabla\cdot F\,dV.
\]
\end{theorem}


Let $\Omega$ be a solid with outward boundary $\partial\Omega$.  For $F=(P,Q,R)$,
\[
  \iint_{\partial\Omega}F\cdot n\,dS
  =
  \iiint_\Omega (P_x+Q_y+R_z)\,dV.
\]
The divergence
\[
  \nabla\cdot F=P_x+Q_y+R_z
\]
measures local source strength.  Gauss' theorem says that total source inside equals outward flux through the boundary.

Use Gauss when the surface is closed, or when it can be closed by adding a simple cap.  If the required surface is not closed, compute the closed flux and subtract the flux across the added part.  If the problem asks for inward orientation, multiply the outward result by $-1$.

\section{Singular fields and missing points}

The note includes the important example type where the field is singular at a point.  Gauss' theorem cannot be applied across a singularity.  The usual remedy is to remove a small sphere around the singular point, apply Gauss' theorem to the punctured region, and then let the radius tend to zero.  For inverse-square radial fields, the flux through the small sphere often remains nonzero.  This is the analytic origin of point charges and delta-type sources.

\section{Stokes theorem}

\begin{theorem}[Stokes theorem]
For an oriented surface \(S\) with boundary \(\partial S\),
\[
  \int_{\partial S} F\cdot dr=\int_S (\nabla\times F)\cdot n\,dS.
\]
\end{theorem}


Let $S$ be an oriented surface with boundary $\partial S$.  Then
\[
  \oint_{\partial S}F\cdot dr
  =
  \iint_S (\nabla\times F)\cdot n\,dS.
\]
Here
\[
  \nabla\times F=
  \begin{vmatrix}
  i&j&k\\
  \partial_x&\partial_y&\partial_z\\
  P&Q&R
  \end{vmatrix}.
\]
The right-hand rule fixes orientation: if the thumb points along $n$, the fingers curl in the positive boundary direction.

Use Stokes when the boundary curve is complicated but some spanning surface is simple, or when the curl is much simpler than the original vector field.  The surface can be changed as long as the boundary and orientation are preserved.


\section{Stokes model computations: area vector and spanning-surface freedom}

\subsection{Triangle in a plane}

Suppose \(\partial S\) is the positively oriented boundary of a triangle with vertices \(A,B,C\), and suppose \(\nabla\times F=c\) is constant.  Stokes gives
\[
  \oint_{\partial S}F\cdot dr
  =\iint_S c\cdot n\,dS
  =c\cdot \vec A_S,
\]
where the oriented area vector is
\[
  \vec A_S=\frac{1}{2}(B-A)\times(C-A).
\]
So the whole line integral collapses to one vector dot product.  The triangle is not merely a convenient region; it packages orientation and area into \(\vec A_S\).

\subsection{Ellipse as a chosen spanning surface}

If the boundary curve is an ellipse lying in a plane and \(\nabla\times F\) is simple, Stokes lets us choose the planar elliptical disk rather than parameterizing the boundary.  If the ellipse has semiaxes \(a,b\) and unit normal \(n\), then its oriented area vector is
\[
  \vec A=\pi ab\,n.
\]
When \(\nabla\times F=c\) is constant,
\[
  \oint_{\partial S}F\cdot dr=c\cdot(\pi ab\,n).
\]
This illustrates the true power of Stokes: the line integral depends on the boundary, but the chosen surface can be any surface spanning it, provided orientation is consistent.

\subsection{Why surface replacement is legal}

If two oriented surfaces \(S_1\) and \(S_2\) have the same boundary, then applying Stokes to the closed surface \(S_1\cup(-S_2)\) gives
\[
  \iint_{S_1}(\nabla\times F)\cdot n\,dS
  -\iint_{S_2}(\nabla\times F)\cdot n\,dS=0,
\]
as long as the field is smooth in the region between them.  If there is a singularity between the two surfaces, this argument can fail.  This is the Stokes-side analogue of the singular-source warning in Gauss' theorem.

\section{Decision table}

\begin{center}
\begin{tabular}{lll}
\toprule
Object & Meaning & Best tool \\
\midrule
$\iint_S f\,dS$ & scalar density on surface & parametrization or graph $dS$ \\
$\iint_SF\cdot n\,dS$ & flux through oriented surface & graph formula or Gauss \\
closed flux & total outward flow & Gauss theorem \\
line circulation & work/circulation along boundary & direct line integral or Stokes \\
\bottomrule
\end{tabular}
\end{center}

\section{Conceptual endpoint}

The high-level statement is the boundary principle:
\[
  \int_{\partial M}\omega=\int_M d\omega.
\]
Green, Gauss, and Stokes are all versions of this formula.  The old vector-calculus notation hides the unity; differential forms reveal it.  The derivative $d\omega$ records local change inside a region, and the boundary integral records the accumulated effect at the edge.


\section{The de Rham complex behind the vector formulas}

The vector identities become less mysterious if the objects are placed by degree:
\[
\begin{array}{c|c|c}
\text{degree} & \text{form in }\mathbb R^3 & \text{vector-calculus name}\\
\hline
0 & f & \text{scalar potential}\\
1 & P\,dx+Q\,dy+R\,dz & \text{work/circulation form}\\
2 & P\,dy\wedge dz+Q\,dz\wedge dx+R\,dx\wedge dy & \text{flux form}\\
3 & g\,dx\wedge dy\wedge dz & \text{volume density}
\end{array}
\]
The exterior derivative moves one row down:
\[
  d:
  \Omega^0\to\Omega^1\to\Omega^2\to\Omega^3.
\]
In vector notation, these three moves are gradient, curl, and divergence.  The identities
\[
  \nabla\times(\nabla f)=0,\qquad
  \nabla\cdot(\nabla\times F)=0
\]
are both the single identity
\[
  d^2=0.
\]
This is why conservative fields, curl-free forms, flux through closed surfaces, and singular sources are connected rather than separate topics.
\section{General Stokes theorem}

\begin{theorem}[General Stokes theorem]
For a compact oriented manifold \(M\) with boundary and a differential form \(\omega\),
\[
  \int_M d\omega=\int_{\partial M}\omega.
\]
\end{theorem}


The vector-calculus theorems in this note are all cases of the general Stokes theorem.  If $M$ is an oriented smooth manifold with boundary and $\omega$ is a compactly supported differential form of degree one less than $\dim M$, then
\[
  \int_{\partial M}\omega=\int_M d\omega.
\]

Green's theorem, Stokes' theorem, and Gauss' theorem correspond to different degrees and dimensions:
\[
  \begin{array}{lll}
  \text{Green} & M\subseteq\mathbb R^2 & \omega=Pdx+Qdy,\\
  \text{Stokes} & M\text{ a surface in }\mathbb R^3 & \omega=Pdx+Qdy+Rdz,\\
  \text{Gauss} & M\subseteq\mathbb R^3 & \omega=Pdy\wedge dz+Qdz\wedge dx+Rdx\wedge dy.
  \end{array}
\]
The apparently different vector formulas are one theorem written in different degrees.

\section{Singular sources and distributions}

Singular fields and missing points lead naturally to distribution theory.  The Dirac delta $\delta_0$ is not an ordinary function but satisfies
\[
  \int \delta_0 \varphi=\varphi(0)
\]
for test functions $\varphi$.

For the inverse-square field in $\mathbb R^3$,
\[
  F(x)=\frac{x}{|x|^3},
\]
one has classically $\nabla\cdot F=0$ away from the origin, but distributionally
\[
  \nabla\cdot F=4\pi \delta_0.
\]
This reconciles the nonzero flux through spheres with zero ordinary divergence away from the singularity.

\section{Conservation laws}

Many PDEs have conservation-law form
\[
  \partial_t \rho+\nabla\cdot J=0.
\]
Integrating over a region $\Omega$ and applying Gauss' theorem gives
\[
  \frac d{dt}\int_\Omega \rho\,dV=-\iint_{\partial\Omega}J\cdot n\,dS.
\]
Change inside equals negative outward flux.  Thus Gauss' theorem is the geometric mechanism behind conservation of mass, charge, and probability.

\section{Currents}

A current generalizes an oriented submanifold by acting on differential forms.  An oriented curve $C$ defines a current
\[
  T_C(\omega)=\int_C \omega.
\]
The boundary of a current is defined by
\[
  \partial T(\omega)=T(d\omega).
\]
This extends Stokes' theorem to singular chains, weak surfaces, and geometric measure theory.

The elementary surface and curve integrals in the note are therefore the smooth case of a much more flexible integration theory over generalized geometric objects.

% Concrete advanced extension
\section{Gauss and Stokes as dual local-to-global laws}

Gauss' theorem and Stokes' theorem are both local-to-global principles.  Gauss says that local source density accumulates to boundary flux.  Stokes says that local rotation accumulates to boundary circulation.  In differential-form notation, both are one formula:
\[
  \int_{\partial M}\omega=\int_M d\omega.
\]
The difference lies only in the degree of the form and the dimension of the manifold.

This unification matters because it prevents memorizing vector identities as separate facts.  The object being integrated determines the theorem: a one-form over a boundary curve gives curl through a surface; a two-form over a closed surface gives divergence through a volume.

\section{Weak forms and PDE}

In PDE, integration by parts and Stokes-type formulas define weak solutions.  Instead of requiring derivatives to exist pointwise, one tests against smooth functions.  For example, a weak divergence equation is read through
\[
  \int_\Omega F\cdot \nabla \varphi\,dV=-\int_\Omega (\nabla\cdot F)\varphi\,dV
\]
up to boundary terms.

Thus vector-calculus integral theorems are not merely computational shortcuts; they are the foundation of weak formulations, conservation laws, and modern PDE theory.
